% !TeX root = ../初等数学笔记.tex
% 由主文档通过 \input 载入；请编译主文档。
\setcounter{chapter}{12} % 保留原书章号。
\chapter{直线、三角形与相似}
\label{chap:em-triangles}
综合几何直接研究点、线、角和长度之间的关系。它的论证始于明确的几何基础，不始于图形“看上去如此”。本章采用欧氏平面的关联、顺序、长度和角度公理及通常的全等基础，不展开完整的公理化体系。

\section{基本对象、角与平行线}
两不同点确定唯一直线；线段包含两端点及其间的点，射线从端点沿一个方向延伸。长度非负并沿同一直线按次序可加，角的大小在指定范围内可加。平角为 $180^\circ$，两直线相交形成的对顶角相等，因为它们都等于平角减去同一邻角。

欧氏平行公设保证过直线外一点恰有一条平行于它的直线。平行线被第三条直线所截，同位角相等、内错角相等、同旁内角互补；反向判定也成立。角的相等关系必须来自这些条件或全等，而非来自线条画得近似平行。过三角形一个顶点作对边的平行线，即可把另外两个内角平移到该顶点附近，三角之和组成平角，故三角形内角和为 $180^\circ$。

\section{三角形的基本性质与全等}
\emph{全等}（congruence）表示两图形可由保持距离的移动完全重合。本书采用下列欧氏三角形基本判定：三边分别相等；两边及夹角分别相等；两角及其夹边分别相等。由内角和可推出“两角及任一对应边”判定，直角三角形还可由斜边与一条直角边判定。对应顶点的次序必须固定，夹角条件不能替换成任意角。

三角形中大边对大角、等边对等角，且任意两边之和大于第三边。后一性质也表达两点之间线段最短：绕经不在该线段上的第三点会增加距离。两边与一个非夹角不足以唯一确定三角形，可能出现两解；第\ref{chap:em-triangle-solving}章将给出其数量条件。退化的三点共线情形不属于本书所谓三角形。

\section{比例、相似与中位线}
两三角形对应角相等、对应边成同一正比例时称为\emph{相似}（similar）。常用判定为两角分别相等、两边成比例且夹角相等、三边成比例。其思路是把其中一个图形按对应边比例缩放到一对边相等，再归结到全等判定；缩放保持角度与长度比例是欧氏相似结构的基础性质。

设 $D\in AB,E\in AC$ 且 $DE\parallel BC$，对应角相等给出 $\triangle ADE\sim\triangle ABC$，所以 $AD/AB=AE/AC=DE/BC$。特别地，连接两边中点的中位线平行于第三边且长为其一半。相似比为 $k>0$ 时，底和高同时乘以 $k$，因而面积乘以 $k^2$。面积比不是一般的长度比，使用时须区分维度。

\section{勾股定理、面积与距离}
\begin{theorem}[勾股定理及逆定理]
直角三角形两直角边长为 $a,b$，斜边长为 $c$，则 $a^2+b^2=c^2$。反之，三角形三边满足此式时，长为 $c$ 的边所对的角为直角。
\end{theorem}
\begin{proof}
在直角顶点向斜边作高，垂足把斜边分成长 $p,q$ 的两段，$p+q=c$。两个小三角形分别与原三角形相似，得到 $a^2=cp,b^2=cq$，相加即得结论。反之，以长 $a,b$ 的两线段构造直角三角形，其斜边由已证公式等于 $c$，与原三角形三边相等，故原对应角也是直角。
\end{proof}
三角形面积为 $S=bh/2$，其中 $b$ 是所选底边长，$h$ 为对该底边所在直线的距离；这可由两份全等三角形拼成平行四边形得到。同高三角形面积比等于底边比，同底时等于高之比。面积法把多条线段关系转化为可加的量，往往比逐个求长更直接。

\section{特殊线段、交点与辅助线}
中线连接顶点与对边中点，角平分线把内角等分，高线垂直于对边所在直线。线段垂直平分线上的点到两端点等距，反向也成立：正向由直角三角形全等，反向由两等腰三角形或垂足两侧的勾股关系得到。

三边垂直平分线共点，因为前两条的交点到三个顶点距离相等，必也在第三条上，该点为外心。两内角平分线交点到三边距离相等，因而位于第三条内角平分线上，得到内心；角平分线等距性质由两个直角三角形全等证明。三条中线共点且把每条中线按从顶点起 $2:1$ 分割：两中线交点与两个边中点构成的三角形，利用中位线与原顶点边平行且为其一半，得到上述比；换另一对中线也得到同一唯一分点。该点称为重心。

三条高线也共点，称为垂心：分别过三个顶点作对边平行线形成外侧三角形，原三角形的顶点是外侧三角形各边中点，原高线就是外侧三角形的垂直平分线，归结到外心共点性。外心、垂心在钝角三角形中可位于外部。若 $AD$ 平分 $\angle A$ 且 $D\in BC$，则 $\triangle ABD,\triangle ACD$ 的面积比一方面为 $BD/DC$，另一方面以 $AB,AC$ 为底，由 $D$ 到两边距离相等得到 $AB/AC$，故 $BD/DC=AB/AC$。

\section{例题、习题与学习检查}
\begin{example}
直角三角形直角边为 $3,4$，求斜边上的高及它分出的两段长度。
\end{example}
\begin{solve}
斜边为 $5$，面积既为 $3\cdot4/2=6$，也为 $5h/2$，故 $h=12/5$。由相似关系，两投影长分别为 $3^2/5=9/5$ 和 $4^2/5=16/5$，其和为 $5$，可作为计算检查。
\end{solve}
\begin{enumerate}
\item 两相似三角形对应边比为 $2:3$，求周长比与面积比。答：$2:3$ 与 $4:9$。
\item 若三角形角平分线把对边分为 $2,3$，其相邻两边之和为 $10$，求两边。答：$4,6$。
\item 证明等腰三角形顶角平分线同时是底边上的中线与高线。提示：比较被分成的两个三角形。
\end{enumerate}

